\documentclass[11pt]{article}

\usepackage[T1]{fontenc}
\usepackage[utf8]{inputenc}
\usepackage{amsmath,amssymb,amsthm}
\usepackage{booktabs}
\usepackage{geometry}
\usepackage{hyperref}
\usepackage{url}

\geometry{margin=1in}

\newtheorem{theorem}{Theorem}
\newtheorem{proposition}{Proposition}
\newtheorem{lemma}{Lemma}
\newtheorem{corollary}{Corollary}
\theoremstyle{definition}
\newtheorem{definition}{Definition}
\theoremstyle{remark}
\newtheorem{remark}{Remark}

\title{Computational Extension of OEIS Sequence A029483:\\
Two New Terms via Block Modular Concatenation in Base 14}
\author{Carlo Corti\\
Independent researcher\\
\texttt{carlo.corti@outlook.com}}
\date{June 2026}

\begin{document}

\maketitle

\begin{abstract}
OEIS sequence A029483 consists of the positive integers \(k\) such that
\(k\) divides the left concatenation of \(1,2,\ldots,k\), written in base
14.  We report a segmented computation over the interval
\[
  [30000000001,300000000000],
\]
finding exactly two new terms in that interval:
\[
  44445214853,\qquad 243987299933.
\]
Together with the fourteen terms already present in the source OEIS record,
these give a proposed b-file extension through \(a(16)\).  The computation
does not construct the concatenated integer.  For each candidate \(k\), the
base-14 left concatenation is evaluated modulo \(k\) by digit-length blocks,
using arithmetic-geometric sums modulo \(k\).  An elementary filter removes
all even candidates and all multiples of \(7\), leaving
\(115714285714\) candidates for the full modular test out of
\(270000000000\) raw integers.  The production search was implemented as a
single C17 program using pthreads and 64-bit candidates with
\texttt{\_\_uint128\_t} modular products.  The two new values, and the
fourteen previously listed values, were independently certified by a Python
checker based on a \(3\times 3\) matrix method, independent of the C
program's recursive series-combination kernel.  The negative part of the
finite search result is based on the validated production kernel and recorded
segment manifest; the independent checker certifies the reported positive
values and known-term replay, not a second exhaustive scan of all non-hits.
\end{abstract}

\noindent\textbf{2020 Mathematics Subject Classification.}
Primary 11Y55; Secondary 11A63, 11B83.

\noindent\textbf{Key words and phrases.}
integer sequence, OEIS, concatenation, base 14, modular arithmetic,
computational number theory, segmented search.

\section{Introduction}

OEIS sequence A029483 is defined as the sequence of positive integers \(k\)
that divide the left concatenation of all integers \(1,2,\ldots,k\), written
in base 14 \cite{oeisA029483}.  In the OEIS source snapshot used for this
project, the known terms were
\[
\begin{gathered}
1,\ 13,\ 143,\ 169,\ 221,\ 403,\ 587,\ 11219,\ 178357,\ 222157,\\
85762339,\ 1086336563,\ 8005332049,\ 15081597011,
\end{gathered}
\]
and the entry recorded that no other terms occur below \(3\cdot 10^{10}\).
The present computation extends the search above that bound to
\(3\cdot 10^{11}\).

A029483 belongs to a broader OEIS family of concatenation-divisibility
sequences in which the base and the direction of concatenation vary.  That
family context is useful provenance, but no transfer theorem between bases or
directions is used here; this paper treats only the base-14 left-concatenation
case.

The word ``left'' is important.  The concatenation considered here is the
base-14 integer obtained by successively prepending the current integer.
Thus, for example, after processing \(1,2,3\), the represented base-14 digit
string is the digit string for \(3\), followed by that for \(2\), followed by
that for \(1\).  Directly constructing this integer is impossible at the
scale considered here: for \(k\) near \(3\cdot 10^{11}\), the concatenation
has trillions of base-14 digits.  The computation instead evaluates the
concatenation modulo the candidate divisor \(k\), independently for each
candidate.

The main result of this paper is the following finite computational
extension.

\begin{theorem}[Validated production search result]
\label{thm:main}
Subject to the recorded segment manifest and the correct execution of the
validated production kernel described below, the only integers \(k\) in the
interval
\[
  [30000000001,300000000000],
\]
satisfying the A029483 divisibility condition are
\[
  44445214853 \quad\text{and}\quad 243987299933.
\]
\end{theorem}

\begin{corollary}
\label{cor:bfile}
Combining Theorem~\ref{thm:main} with the source OEIS record used as input
to this project, the proposed extension of the b-file for A029483 through the
searched range has sixteen terms, with
\[
  a(15)=44445214853,\qquad a(16)=243987299933.
\]
\end{corollary}

The rest of the paper explains the computation supporting
Theorem~\ref{thm:main}.  Section~\ref{sec:def} fixes notation for left
concatenation.  Section~\ref{sec:block} gives the block modular evaluation
used by the production program.  Section~\ref{sec:search} describes the
filter and segmented search, while Section~\ref{sec:specificity} records the
particular computational interpretation of this search.  Section~\ref{sec:impl}
records implementation and reproducibility details, and
Section~\ref{sec:validation} records the independent certification and
validation evidence.
Section~\ref{sec:results} summarizes the new terms and the proposed b-file
extension.

\subsection*{Notation}

For an integer \(b\ge 2\) and \(n\ge 1\), let
\[
  \ell_b(n)=1+\lfloor \log_b n\rfloor
\]
be the number of base-\(b\) digits of \(n\).  Let
\[
  L_b(n)=\sum_{j=1}^{n}\ell_b(j), \qquad L_b(0)=0.
\]
For later use, if \(d\ge 1\), then
\[
  L_b(b^{d-1}-1)=\sum_{i=1}^{d-1} i(b^i-b^{i-1}),
\]
with the empty sum equal to \(0\).  This is the total number of base-\(b\)
digits contributed by the positive integers with fewer than \(d\) digits.
All intervals of integers in this paper are inclusive unless explicitly
stated otherwise.

\section{Definition of the left concatenation}
\label{sec:def}

\begin{definition}
For \(b\ge 2\) and \(k\ge 1\), define the left concatenation
\[
  C_b(k)=\sum_{n=1}^{k} n\,b^{L_b(n-1)}.
\]
The sequence considered here is
\[
  A029483=\{k\ge 1 : C_{14}(k)\equiv 0 \pmod{k}\}.
\]
\end{definition}

This formula agrees with the OEIS Mathematica program in the local source
snapshot, in which the base-14 digits of the current integer are prepended to
the digit list already constructed.  The term \(n\) is shifted left by exactly
the number of base-14 digits already present, namely \(L_b(n-1)\).

\section{Block modular evaluation}
\label{sec:block}

The computation evaluates \(C_{14}(k)\bmod k\) without constructing
\(C_{14}(k)\).  The key observation is that the summation naturally splits
into blocks on which all integers have the same base-14 digit length.

\begin{proposition}[Digit-length block formula]
\label{prop:block}
Let \(b=14\).  Fix \(k\ge 1\).  For a digit length \(d\ge 1\), set
\[
  a=b^{d-1}, \qquad u=\min(k,b^d-1), \qquad q=b^d,
\]
and suppose \(a\le k\).  Let \(N=u-a+1\) and \(s=L_b(a-1)\).  The
contribution to \(C_b(k)\) from the block \(a\le n\le u\) is
\[
  b^s\sum_{r=0}^{N-1}(a+r)q^r \pmod{k}.
\]
Consequently,
\[
  C_b(k)\equiv
  \sum_{\substack{d\ge 1\\ b^{d-1}\le k}}
  b^{L_b(b^{d-1}-1)}
  \sum_{r=0}^{\min(k,b^d-1)-b^{d-1}}(b^{d-1}+r)(b^d)^r
  \pmod{k}.
\]
\end{proposition}

\begin{proof}
For \(n=a+r\) in the digit-length \(d\) block, the integers preceding \(n\)
in the left-concatenation construction contribute
\[
  L_b(n-1)=L_b(a-1)+dr=s+dr
\]
base-\(b\) digits.  Hence the summand in \(C_b(k)\) is
\[
  (a+r)b^{s+dr}=b^s(a+r)(b^d)^r.
\]
Summing over \(0\le r<N\) gives the block contribution.  Summing the
contributions over all digit lengths up to \(k\) gives the stated congruence.
\end{proof}

For a fixed \(q\), \(N\), and modulus \(k\), the inner sum can be written
\[
  aS_0(q,N)+S_1(q,N),
\]
where
\[
  S_0(q,N)=\sum_{r=0}^{N-1}q^r,\qquad
  S_1(q,N)=\sum_{r=0}^{N-1}r q^r.
\]
The production C program computes the triple
\[
  \bigl(q^N,S_0(q,N),S_1(q,N)\bigr)\pmod{k}
\]
by recursive series composition.  The recursion splits a requested length
\(N\) by halving it, or by removing a final odd term and then halving, and
combines the resulting subsegments; it is not a linear scan through all
\(N\) summands.  In this computational discussion \(q=b^d\), with its
residue modulo \(k\) used in the modular arithmetic.  The base cases are
\[
  (q^0,S_0(q,0),S_1(q,0))=(1,0,0),\qquad
  (q^1,S_0(q,1),S_1(q,1))=(q,1,0).
\]
If a left segment has length \(m\) and a
right segment has length \(n\), then
\[
\begin{aligned}
S_0(q,m+n)
  &= S_0(q,m)+q^mS_0(q,n),\\
S_1(q,m+n)
  &= S_1(q,m)+q^m\bigl(mS_0(q,n)+S_1(q,n)\bigr).
\end{aligned}
\]
Here \(m\) denotes the length of the left subsegment; in modular
implementation it is reduced modulo \(k\) before being multiplied by a
residue.
These identities allow binary splitting in \(N\), so each digit-length block
is evaluated by a logarithmic number of modular products.  Termination follows
because each recursive call strictly reduces the requested segment length
until one of the two displayed base cases is reached; correctness then follows
by induction from the combination identities.  This is the standard
divide-and-conquer idea behind binary splitting
\cite{brentZimmermann2010}.  For a fixed candidate \(k\), the number of
digit-length blocks is \(O(\log k)\) and each block contributes
\(O(\log k)\) modular products, giving a qualitative
\(O((\log k)^2)\) modular-product cost per candidate for fixed base.

\begin{remark}
No division by \(q-1\) is used.  This avoids special cases when \(q-1\) is not
invertible modulo \(k\).
\end{remark}

\section{Filtering and segmented search}
\label{sec:search}

The production search used one elementary but exact filter.

\begin{lemma}[Base-factor exclusion]
\label{lem:filter}
For \(k>1\), if \(k\) is even or divisible by \(7\), then \(k\notin A029483\).
\end{lemma}

\begin{proof}
For \(b=14\), every summand \(n b^{L_b(n-1)}\) with \(n\ge 2\) is divisible
by both \(2\) and \(7\), since \(b=14=2\cdot 7\) and \(L_b(n-1)>0\).  The
\(n=1\) summand is \(1\).
Therefore
\[
  C_{14}(k)\equiv 1\pmod 2,\qquad C_{14}(k)\equiv 1\pmod 7.
\]
If \(k>1\) were even and divided \(C_{14}(k)\), then \(C_{14}(k)\) would be
congruent to \(0\) modulo \(2\), a contradiction.  The argument modulo \(7\)
is identical.
\end{proof}

The search interval was
\[
  [30000000001,300000000000],
\]
containing \(270000000000\) raw integers.  Lemma~\ref{lem:filter} leaves
\(115714285714\) candidates for the full modular test.  The production run
was split into 23 consecutive segments, normally of \(12000000000\) raw
integers each, with the last segment shortened to end exactly at
\(300000000000\).  The exact segment record is stored in
\path{results/A029483_official_run_segments.tsv} in the raw package
\cite{cortiA029483Package2026}.
The base-14 digit-length boundary \(14^{10}=289254654976\) lies inside the
searched interval; candidates above this boundary are handled automatically by
the same digit-length block formula.
The survivor count can be checked directly from the six admissible residue
classes modulo \(14\):
\[
  \sum_{r\in\{1,3,5,9,11,13\}}
  \left(
    \left\lfloor\frac{300000000000-r}{14}\right\rfloor
    -
    \left\lfloor\frac{30000000000-r}{14}\right\rfloor
  \right)
  =115714285714.
\]
Table~\ref{tab:aggregate} summarizes the aggregate run, and
Table~\ref{tab:hitsegments} lists the two segments in which new terms were
found.

\begin{table}[ht]
\centering
\begin{tabular}{rrrr}
\toprule
Raw integers & Survivors after filter & Segments & New hits\\
\midrule
270000000000 & 115714285714 & 23 & 2\\
\bottomrule
\end{tabular}
\caption{Aggregate production search over
\([30000000001,300000000000]\).}
\label{tab:aggregate}
\end{table}

\begin{table}[ht]
\centering
\begin{tabular}{rrr}
\toprule
Segment interval & Hit & Cumulative new hits\\
\midrule
\([42000000001,54000000000]\) & 44445214853 & 1\\
\([234000000001,246000000000]\) & 243987299933 & 2\\
\bottomrule
\end{tabular}
\caption{The two production segments containing new A029483 terms.}
\label{tab:hitsegments}
\end{table}

The wall-clock time recorded across the segment manifest is
\(70941.341\) seconds, or approximately \(19.706\) hours.  This timing is
reported as a property of the recorded production run, not as a hardware
independent benchmark.

\section{What is specific to this search}
\label{sec:specificity}

The present computation is best understood as a certified divisibility search,
not as a statistical experiment.  In contrast with prime-counting or
square-value problems, A029483 does not come with an evident density model
whose predictions can be checked numerically in a short paper.  The divisor
being tested is also the modulus in which the concatenation is evaluated, and
the digit-length boundaries of the base-14 expansion enter the predicate
directly.  This makes the search highly deterministic but does not by itself
suggest a credible random model for the location of the next term.

The one simple structural filter is nevertheless exact and unusually cheap:
because the rightmost digit of the left concatenation is always the digit of
\(1\), every candidate divisible by a prime factor of the base is impossible
for \(k>1\).  For base \(14\), this removes all even candidates and all
multiples of \(7\).  Over a complete residue period modulo \(14\), this leaves
\(3/7\) of the integers; in the present finite interval the exact survivor
count is \(115714285714\).  The remaining work is not a sieve cascade but a
repeated exact modular evaluation of \(C_{14}(k)\) by digit-length blocks.

For this reason, the two new terms should not be read as evidence for a
particular asymptotic frequency.  Their significance is more modest and more
concrete: they extend the certified table beyond the previous OEIS-recorded
bound, and the segment manifest gives an explicit negative result for every
other integer in the same interval.

\section{Implementation and reproducibility}
\label{sec:impl}

The production program is the C17 source file \path{src/a029483.c}, built
with the accompanying \path{src/Makefile}.  The implementation targets the C17
language standard \cite{isoC17} and uses POSIX threads \cite{posix2017}.
Candidates and residues are stored in unsigned 64-bit integers.  Modular
products use the GCC/Clang extension \path{__uint128_t} as a double-width
intermediate before reduction modulo the current candidate \(k\).  The program
evaluates each candidate independently by the block formula of
Proposition~\ref{prop:block}, so segmenting and resuming do not require
carrying a concatenation state from one candidate modulus to the next.
The largest modulus tested here is at most \(3\cdot 10^{11}\); before
reduction, a product of two residues is therefore at most \((k-1)^2\), and
\[
  (k-1)^2 < k^2 \le 9\cdot 10^{22} < 2^{77},
\]
well within the \(128\)-bit intermediate range.  In the binary-splitting
combination for \(S_1\), the segment length \(m\) is first reduced modulo
\(k\); consequently every multiplication passed to the double-width
intermediate is still a product of two residues modulo \(k\).

The production and validation host recorded for this project was a MinisForum
UM790 Pro with an AMD Ryzen 9 7940HS processor, 8 physical cores and 16
hardware threads, and 64 GB of DDR5 memory.  The local software environment
used for the package validation was Linux Mint 22.3, Linux kernel
6.14.0-37-generic, and GCC 13.3.0.  The Makefile mandatory flags were
\begin{verbatim}
-std=c17 -Wall -Wextra -Wshadow -Wconversion
-Wdouble-promotion -Wformat=2 -pthread
\end{verbatim}
with the optimized production build adding
\begin{verbatim}
-O3 -march=native -mtune=native -flto
-funroll-loops -fomit-frame-pointer
\end{verbatim}
and linking with \texttt{-flto -pthread -lm}.

The official command pattern was
\begin{verbatim}
src/a029483 \
  --resume results/a029483_checkpoint.txt \
  --end END \
  --block-size 12000000000 \
  --threads 16 \
  --out results/A029483_new_terms_u64.tsv \
  --checkpoint results/a029483_checkpoint.txt \
  --progress-every 1
\end{verbatim}
with \texttt{END} set by segment.  The \texttt{--block-size} argument is a
raw-integer span before filtering, not the number of surviving candidates.
The \texttt{--resume} option reads the next start from an existing checkpoint;
\texttt{--checkpoint} names the checkpoint file to update after a completed
segment.  This command is shown as run from the extracted raw package root.
The public raw package validation summary records that the following checks,
shown from the parent project directory containing \path{GitHub/}, passed:
\begin{verbatim}
make -C GitHub/src -f Makefile clean
make -C GitHub/src -f Makefile test
make -C GitHub/src -f Makefile sanitize
GitHub/scripts/certify_a029483.py --known --new
GitHub/scripts/certify_a029483.py --file GitHub/results/A029483_new_terms_u64.tsv
\end{verbatim}
The \texttt{test} target runs the internal self-test, a small-range regression
over \([1,250000]\) against the expected known terms in that interval, command
line validation, and checkpoint error checks.  The \texttt{sanitize} target
reruns representative checks under ASan/UBSan.  The two certifier invocations
above replay both the previously known terms and the two new terms, and then
verify the production hit file.

Segments are inclusive.  Each invocation advances from the checkpoint file and
sets a segment endpoint by \texttt{--end}; the checkpoint records the last
completed endpoint, the next start, and the hit count, so a resumed run does
not silently skip or duplicate an interval.  A checkpoint is written only
after the corresponding segment has completed successfully; it is written to a
temporary file, flushed, synced, and then promoted to the checkpoint path.
Candidate tests are independent.  The parallel implementation assigns
candidate ranges to worker threads and serializes shared hit, counter, and
checkpoint updates.  The validation record reports cross-thread determinism on
\([1,250000]\) for thread counts \(1,2,4,8,16\), as a regression check of the
reported term set rather than a byte-for-byte claim about timing-bearing
checkpoint files.

Before the production campaign, the final code had also passed a release
build with strict warnings, the Makefile regression suite, ASan/UBSan, TSan,
cross-thread determinism checks for thread counts \(1,2,4,8,16\), Valgrind
self-test and smoke runs, scan-build, cppcheck triage, clang-tidy triage, and
a throughput probe near \(3\cdot 10^{10}\).  These records are summarized in
\path{validation/validation_summary.md} and
\path{validation/final_code_validation.md}.

\section{Independent certification}
\label{sec:validation}

The independent checker \path{scripts/certify_a029483.py} intentionally does
not reuse the C program's recursive arithmetic-geometric series combination.
It uses the same digit-length block decomposition proved in
Proposition~\ref{prop:block}, but evaluates each block through a linear
recurrence represented by a \(3\times 3\) matrix transition.  The independence is
therefore kernel-level independence,
not a second derivation of the block decomposition.  The checker is a Python 3
script using only the standard library.

For a block beginning at \(a\), with \(q=b^d\bmod k\), define the state
\[
  v_i=\begin{pmatrix}q^i\\ i q^i\\
  \sum_{r=0}^{i-1}(a+r)q^r\end{pmatrix}\pmod{k}.
\]
Then
\[
  v_{i+1} =
  \begin{pmatrix}
  q & 0 & 0\\
  q & q & 0\\
  a & 1 & 1
  \end{pmatrix}
  v_i \pmod{k},
  \qquad
  v_0=\begin{pmatrix}1\\0\\0\end{pmatrix}.
\]
Raising this \(3\times 3\) matrix to the block length gives the same block
sum as Proposition~\ref{prop:block}, but through a different computational
kernel: the third component advances by adding
\((a+i)q^i\) to the accumulated block sum.  This provides an independent
certification path for reported values.

The validation summary records
\begin{verbatim}
k=44445214853: residue=0 direct_residue=n/a PASS
k=243987299933: residue=0 direct_residue=n/a PASS
\end{verbatim}
for the two new values.  The checker also certified the fourteen terms from
the source OEIS record.  In these large cases, \texttt{direct\_residue=n/a}
means that the checker did not attempt literal construction of the
concatenated integer; that direct path is reserved for small values where the
integer is feasible to build.  With its default setting, the checker compares
the matrix residue against literal left-concatenation residues for
\(k\le 5000\), which anchors the shared block decomposition on small cases.
The reported residue is still computed modulo \(k\) by the independent matrix
method above.  During the campaign, a
neighbor check around the second new term recorded
\begin{verbatim}
k=243987299932: residue=196861283305 direct_residue=n/a FAIL
k=243987299933: residue=0 direct_residue=n/a PASS
k=243987299934: residue=150270907739 direct_residue=n/a FAIL
\end{verbatim}
where \texttt{FAIL} denotes the expected non-zero residue for a non-term, not
a program failure.  This is a local sanity check on the reported hit; the two
neighboring residues are non-zero modulo their respective candidates.

\section{Results}
\label{sec:results}

The new terms found by the production search are stored in
\texttt{results/A029483\_new\_terms\_u64.tsv}:
\[
  44445214853,\qquad 243987299933.
\]
The proposed b-file extension stored in \texttt{data/b029483.txt} is shown in
Table~\ref{tab:bfile}.  The companion file
\texttt{data/certified\_terms.tsv} records source and certificate status for
all sixteen rows; it is a certification table, not an OEIS a-file.

\begin{table}[ht]
\centering
\begin{tabular}{rr}
\toprule
\(n\) & \(a(n)\)\\
\midrule
1 & 1\\
2 & 13\\
3 & 143\\
4 & 169\\
5 & 221\\
6 & 403\\
7 & 587\\
8 & 11219\\
9 & 178357\\
10 & 222157\\
11 & 85762339\\
12 & 1086336563\\
13 & 8005332049\\
14 & 15081597011\\
15 & 44445214853\\
16 & 243987299933\\
\bottomrule
\end{tabular}
\caption{Proposed b-file extension for A029483 through the two terms found
in this project.}
\label{tab:bfile}
\end{table}

\begin{proof}[Proof of Theorem~\ref{thm:main}]
This is a computer-assisted finite proof.  The official segment manifest
partitions the interval \([30000000001,300000000000]\) into 23 consecutive
inclusive segments.  The raw counts in the manifest sum to \(270000000000\),
the size of the interval.
By Lemma~\ref{lem:filter}, every even integer and every multiple of \(7\) in
the interval is impossible.  The same manifest records
\(115714285714\) surviving candidates after this exact filter, and these are
the candidates evaluated by the block modular test of
Proposition~\ref{prop:block}.

The production hit file \texttt{results/A029483\_new\_terms\_u64.tsv}
contains exactly the two rows \(44445214853\) and \(243987299933\).  Both
values lie in the interval and both are certified independently by the
matrix-based checker described in Section~\ref{sec:validation}.  The negative
claim depends on the integrity of the recorded segment manifest and the
correct execution of the validated production kernel; the independent
certificate confirms the reported positive hits and known-term replay, but it
is not a second exhaustive scan of every non-hit.  Subject to those recorded
computational artifacts, the non-surviving candidates are excluded by
Lemma~\ref{lem:filter}, and every surviving candidate in the segment
partition was evaluated by the production program.  Hence the two certified
values are exactly the A029483 terms found in the searched interval.
\end{proof}

\begin{proof}[Proof of Corollary~\ref{cor:bfile}]
The local OEIS source snapshot lists the first fourteen terms and records no
other terms below \(3\cdot 10^{10}\).  This earlier negative bound is taken as
input from that OEIS snapshot and is not re-derived in the present search.
The present computation begins at \(30000000001\), the first integer strictly
above the previously recorded bound.  Theorem~\ref{thm:main} then accounts for
the interval from \(30000000001\) through \(300000000000\), adding exactly the
two terms shown in Table~\ref{tab:bfile}.  The displayed indices \(a(15)\) and
\(a(16)\) therefore depend on the local OEIS source snapshot, including its
negative statement below \(3\cdot 10^{10}\).  A live check of the OEIS entry on
2026-06-03 showed the same fourteen terms and the same stated bound.
\end{proof}

\section{Conclusion}

The search reported here changes the known data for A029483 immediately above
the previously recorded bound.  In the interval from \(30000000001\) through
\(300000000000\), exactly two terms were found:
\[
  44445214853,\qquad 243987299933.
\]
Consequently the proposed b-file extension has sixteen terms, with these two
values occupying \(a(15)\) and \(a(16)\).  The computational content of the
paper is therefore a finite certified extension of the table, not a statistical
or heuristic model for the occurrence of such divisors.

No density law for A029483 is proposed here, and no probabilistic calculation
is used to support the result.  This is a deliberate limitation rather than an
omission: for this sequence the reliable contribution is the certified
extension of the table and the explicit negative result over the searched
interval.  The exclusion of all other integers in that interval rests only on
the exact base-factor filter, the exhaustive block modular evaluation of the
surviving candidates, and the independent matrix-based certification of the
reported terms.  A final live OEIS check was performed on 2026-06-03 and
confirmed that the source entry still listed the same fourteen terms and the
same bound below \(3\cdot 10^{10}\).

\section*{Data availability}

The code, validation records, segment manifest, production hit file, proposed
b-file, and certified-term table are part of the local project package and are
to be included in the project repository and archive listed in the References
when those public artifacts are finalized.  The principal artifacts are the C17
source and Makefile, the official segment TSV, the new-term TSV, the proposed
b-file, the certified-term TSV, the independent Python checker, and the
validation summaries.  No additional dataset beyond these project artifacts is
needed to reproduce the reported finite extension.

\section*{AI use disclosure}

This work used large language models as research assistants during the
workflow for the A029483 project.  DeepSeek was used in the initial triage of
the sequence and project feasibility.  An OpenAI Codex system labelled GPT-5
by the working environment was used during the code-review,
manuscript-drafting, and manuscript-revision stages.  Additional LLM systems,
identified as Claude, Gemini, Grok, Kimi, Qwen, and Z.ai, were used as
non-author peer reviewers of draft manuscript text.

The AI-assisted work included mathematical analysis of the base-14
left-concatenation predicate, design and refinement of the digit-length block
evaluation, C17 implementation support, code-review triage, validation
planning, and drafting of this manuscript.

Mathematical analysis.  The modular block formula, the arithmetic-geometric
series presentation, and the elementary exclusion of even candidates and
multiples of \(7\) were developed and written with AI assistance.  The
identities used in the paper are proved explicitly in the text and were
checked against the final implementation and the independent certificate
script before inclusion.

Algorithm and code.  The production C17 search program and Makefile were
produced through an AI-assisted development and review loop.  Their correctness
does not rest on that provenance: it rests on the regression suite, sanitizer
and tooling checks, cross-thread determinism checks, known-term replay,
segmented production manifest, and the independent Python matrix certificate
described in this paper.

Manuscript.  The manuscript was drafted with AI assistance and revised by the
author.  Every quantitative claim included here was checked against the local
computational artifacts: the OEIS source snapshot, the segment manifest, the
new-term file, the proposed b-file, the certified-term table, and the
validation summaries.

The author takes full scientific responsibility for the mathematical
statements, the algorithm, the implementation, the numerical results, and the
conclusions of this paper.

\begin{thebibliography}{9}
\sloppy
\raggedright

\bibitem{oeisA029483}
OEIS Foundation Inc.,
\emph{Sequence A029483: Numbers \(k\) that divide the (left)
concatenation of all numbers \(\le k\) written in base 14},
The On-Line Encyclopedia of Integer Sequences,
\url{https://oeis.org/A029483}.
Sequence created by O. G\'erard; further terms by L. Reeves, Jul. 09, 2001;
edited and updated by L. Reeves, Apr. 12, 2002; Mathematica program by
R. Price, Mar. 12, 2020; \(a(11)\) contributed by M. Alekseyev, May 15, 2011;
\(a(12)\)--\(a(14)\) contributed by J. Yuen, Jun. 04, 2024.  Local project
snapshot used as source material, 2026-05-29; live check performed
2026-06-03.

\bibitem{isoC17}
International Organization for Standardization,
\emph{ISO/IEC 9899:2018, Information technology -- Programming languages -- C},
International Organization for Standardization, Geneva, 2018.

\bibitem{posix2017}
IEEE and The Open Group,
\emph{The Open Group Base Specifications Issue 7, 2018 edition; IEEE Std
1003.1-2017},
2018.  Available at \url{https://pubs.opengroup.org/onlinepubs/9699919799/}.

\bibitem{brentZimmermann2010}
R. P. Brent and P. Zimmermann,
\emph{Modern Computer Arithmetic},
Cambridge University Press, Cambridge, 2010.

\bibitem{cortiA029483Package2026}
C. Corti,
\emph{A029483: Computational search package for OEIS sequence A029483},
GitHub repository and Zenodo archive, 2026.
\url{https://github.com/carcorti/A029483}.
\url{https://doi.org/10.5281/zenodo.20538439}.

\end{thebibliography}

\end{document}
